\documentclass[11pt,a4paper]{article}

% ---------------------------------------------------------
% PACKAGES
% ---------------------------------------------------------
\usepackage[utf8]{inputenc}
\usepackage[T1]{fontenc}
\usepackage{lmodern}
\usepackage{amsmath, amssymb, amsfonts, bm}
\usepackage{mathtools}
\usepackage{physics}
\usepackage{tensor}
\usepackage{bbm}
\usepackage{graphicx}
\usepackage{float}
\usepackage{cite}
\usepackage{geometry}
\geometry{margin=2.5cm}
\usepackage{microtype}
\usepackage[most]{tcolorbox}

% ---------------------------------------------------------
% HYPERREF
% ---------------------------------------------------------
\usepackage[colorlinks=true,
            linkcolor=blue,
            citecolor=blue,
            urlcolor=blue,
            pdfauthor={Michael Lehmann},
            pdftitle={Observational Synthesis and Global Fit in RDCC},
            pdfsubject={Relaxation-Driven Cyclic Cosmology},
            pdfkeywords={RDCC, cosmology, global fit, CPT symmetry, relaxation dynamics}
           ]{hyperref}

% ---------------------------------------------------------
% TITLE
% ---------------------------------------------------------
\title{\textbf{Dynamical Generation of the Relaxation Parameter in\\
Relaxation-Driven Cyclic Cosmology (RDCC)}}

\author{Michael Lehmann\\
Independent Researcher\\
\texttt{mi.lehmann@gmx.de}}

\date{April 2026}

\begin{document}

\maketitle

\begin{abstract}
Relaxation-Driven Cyclic Cosmology (RDCC) is a one-parameter, CPT-symmetric
framework in which all observable deviations from $\Lambda$CDM arise from the
infrared value of the relaxation parameter $\alpha_{\mathrm{IR}}$. While Companions~I--V
established the gravitational foundations, tensor phenomenology, relaxation-induced
scalar sector, global CPT state, and the first prototype global fit, the microscopic
time evolution of the relaxation parameter has not yet been derived on a realistic
cosmological background.

This paper, Companion~VI, closes this gap. We derive the relaxation equation
$\dot{\alpha} = \Gamma \alpha$ from the dimension-six inter-sector operator, show that the
naive EFT estimate $\Gamma \sim H^5/M^4$ fails to generate any observable relaxation,
and introduce the minimal generalized rate $\Gamma = \beta H$ consistent with the
CPT-symmetric structure. Matching to the observed value $\alpha_{\mathrm{IR}} \simeq 0.0175$
fixes $\beta \simeq 1.25 \times 10^{-4}$.

We solve the full evolution equation $\mathrm{d}\alpha/\mathrm{d}a = (\Gamma/H)\alpha$ on a
$\Lambda$CDM background and find that $\alpha(a)$ grows sharply near the bounce and
freezes to its infrared value long before matter–radiation equality. This dynamical
behavior explains why the RDCC prediction vector remains stable across more than ten
orders of magnitude in scale and provides a natural microscopic origin for the observed
value of $\alpha_{\mathrm{IR}}$.

Companion~VI thereby completes the dynamical foundation of RDCC and prepares the
ground for a full Boltzmann implementation in Companion~VII.
\end{abstract}

\newpage

\tableofcontents

\newpage

\section{Introduction}

Relaxation-Driven Cyclic Cosmology (RDCC) describes the Universe as a globally coherent,
CPT-symmetric bipartite quantum system whose observable deviations from $\Lambda$CDM are
governed by a single infrared quantity: the relaxation parameter $\alpha_{\mathrm{IR}}$. In the global
Hilbert-space picture developed in Companions~I–IV, this parameter measures the residual loss
of coherence between the visible and mirror sectors after the bounce. Its infrared value controls
the entire RDCC prediction vector,


\[
\mathcal{O}_{\mathrm{RDCC}} = (n_s,\ \Delta N_{\mathrm{eff}},\ f\sigma_8,\ \Omega_{\mathrm{GW}}),
\]


and determines the correlated deviations from $\Lambda$CDM across more than ten orders of
magnitude in scale.

In Companions~I–V, $\alpha_{\mathrm{IR}}$ was treated as an effectively constant infrared quantity:
it vanishes at the CPT-symmetric bounce, grows as the Universe departs from the global
fixpoint, and freezes to its late-time value long before any observable epoch. This assumption
underlies the one-parameter structure of RDCC and the global fit developed in Companion~V.
However, the microscopic time evolution of $\alpha(a)$ has not yet been derived on a realistic
cosmological background. This missing step is essential for completing the internal consistency
of RDCC: the predictive structure of the theory requires that the relaxation parameter freeze
early enough for all observables to depend on the same infrared value.

The purpose of this paper, Companion~VI, is to derive and solve the relaxation equation
governing $\alpha(a)$, evaluate the naive EFT estimate for the relaxation rate, introduce the
minimal generalized rate consistent with the CPT-symmetric structure, and compute the full
numerical evolution of $\alpha(a)$ from the bounce to today. We show that the naive estimate
$\Gamma \sim H^5/M^4$ fails by many orders of magnitude, producing no observable relaxation.
We then introduce the minimal curvature-sensitive generalization $\Gamma = \beta H$, derive the
effective coupling $\beta$ from the observed value $\alpha_{\mathrm{IR}} \simeq 0.0175$, and solve the
full evolution equation on a $\Lambda$CDM background.

The resulting picture is simple and robust: $\alpha(a)$ grows rapidly as the Universe departs
from the CPT-symmetric fixpoint, freezes by $a \sim 10^{-4}$, and remains constant thereafter.
This early freezing explains why the RDCC prediction vector is stable across cosmic history and
provides a natural microscopic origin for the observed value of $\alpha_{\mathrm{IR}}$. Companion~VI
thereby completes the dynamical foundation of RDCC and prepares the ground for a full
Boltzmann implementation in Companion~VII.

\section{The Relaxation Equation}

The relaxation parameter $\alpha$ originates from the unique dimension-six inter-sector operator
introduced in Companions~I and~III,
\begin{equation}
    \mathcal{L}_{\mathrm{int}}
    = \frac{1}{M^2 + \mu^4}\,
      T^{(+)}_{\mu\nu} T^{(-)\,\mu\nu},
\end{equation}
which induces a slow exchange of energy, entropy, and information between the visible and
mirror sectors. In the global Hilbert-space picture of Companion~IV, this operator governs the
rate at which the two CPT-conjugate branches of the Universe decohere as the system departs
from the bounce. The purpose of this section is to derive the relaxation equation governing
$\alpha(a)$ and express it in a form suitable for cosmological evolution on a $\Lambda$CDM
background.

\subsection{Derivation from the Inter-Sector Operator}

Following the analysis of Companion~III, the interaction modifies the continuity equations of the
two sectors,
\begin{align}
    \dot{\rho}_+ + 3H(\rho_+ + p_+) &= +Q, \\
    \dot{\rho}_- + 3H(\rho_- + p_-) &= -Q,
\end{align}
where $Q$ is the inter-sector energy-exchange term. To leading order in the EFT expansion,
\begin{equation}
    Q \sim \frac{H^2}{M^2}\,\Delta T,
\end{equation}
with $\Delta T$ the difference between the sectoral stress-energy traces. The relaxation parameter
is defined as the dimensionless measure of inter-sector equilibration,
\begin{equation}
    \alpha = \frac{\Delta \rho}{\rho_{\mathrm{tot}}},
\end{equation}
and its evolution follows from differentiating this definition and inserting the modified continuity
equations. To leading order in $\alpha$ (the regime relevant for RDCC), the evolution reduces to
the multiplicative form
\begin{equation}
    \dot{\alpha} = \Gamma\,\alpha,
\end{equation}
where $\Gamma$ is the effective relaxation rate determined by the microscopic interaction
strength. This equation expresses the fact that the relaxation dynamics correspond to a
multiplicative decoherence process in the bipartite Hilbert space.

\subsection{Changing Variables to the Scale Factor}

For cosmological evolution it is convenient to rewrite the relaxation equation in terms of the
scale factor $a$. Using
\begin{equation}
    \frac{d}{dt} = aH\,\frac{d}{da},
\end{equation}
the relaxation equation becomes
\begin{equation}
    \frac{d\alpha}{da}
    = \frac{\Gamma}{aH}\,\alpha.
\end{equation}
This is the master equation governing the evolution of the relaxation parameter. Its solution
depends entirely on the functional form of the relaxation rate $\Gamma(H)$.

The next section evaluates the naive EFT estimate $\Gamma \sim H^5/M^4$ and shows that it
fails to generate any observable relaxation over cosmic history, motivating the generalized rate
introduced in Section~4.

\section{Naive EFT Estimate and Its Failure}

The relaxation rate $\Gamma$ appearing in the master equation
\begin{equation}
    \frac{d\alpha}{da}
    = \frac{\Gamma}{aH}\,\alpha
\end{equation}
is determined by the microscopic structure of the dimension-six inter-sector operator. A natural
first estimate follows from dimensional analysis: since the only dynamical scale in the late
Universe is the Hubble parameter $H$, the naive EFT expectation is
\begin{equation}
    \Gamma_{\mathrm{naive}} \sim \frac{H^5}{M^4}.
\end{equation}
This section evaluates this estimate on a realistic $\Lambda$CDM background and shows that it
fails by an enormous margin. The failure is not a numerical subtlety but a structural consequence
of the extreme suppression of higher-dimensional operators in a low-curvature Universe. This
motivates the generalized rate introduced in Section~4.

\subsection{Master Equation with the Naive Rate}

Inserting the naive estimate into the master equation yields
\begin{equation}
    \frac{d\alpha}{da}
    = \frac{H^5}{aH\,M^4}\,\alpha
    = \frac{H^4}{aM^4}\,\alpha.
\end{equation}
Integrating from the bounce ($a_{\min}$) to today ($a_0 = 1$) gives
\begin{equation}
    \ln\!\left(\frac{\alpha_0}{\alpha_{\min}}\right)
    = \int_{a_{\min}}^{1}
      \frac{H(a)^4}{aM^4}\,da.
\end{equation}
Because $H(a)$ is largest at the earliest times, the integral is dominated by the radiation-dominated
regime.

\subsection{Evaluation on a \texorpdfstring{$\Lambda$CDM}{Lambda CDM} Background}

On a flat $\Lambda$CDM background,
\begin{equation}
    H^2(a) = H_0^2\left(\Omega_r a^{-4} + \Omega_m a^{-3} + \Omega_\Lambda\right).
\end{equation}
In the radiation-dominated regime ($a \ll 10^{-4}$),
\begin{equation}
    H^4(a) \simeq H_0^4\,\Omega_r^2\,a^{-8}.
\end{equation}
Substituting into the integral,
\begin{equation}
    \ln\!\left(\frac{\alpha_0}{\alpha_{\min}}\right)
    \sim \frac{H_0^4\,\Omega_r^2}{M^4}
         \int_{a_{\min}}^{a_{\mathrm{rad}}} a^{-9}\,da
    \sim \frac{H_0^4\,\Omega_r^2}{M^4}\,a_{\min}^{-8}.
\end{equation}
Using standard cosmological parameters and a conservative bounce scale
$a_{\min} \sim 10^{-30}$, we obtain
\begin{equation}
    \ln\!\left(\frac{\alpha_0}{\alpha_{\min}}\right)
    \sim 10^{-200}
    \qquad (M \sim 10^{16}\,\mathrm{GeV}).
\end{equation}
Thus,
\begin{equation}
    \alpha(a) \simeq \alpha_{\min}
    \quad\text{for all } a.
\end{equation}

\subsection{Interpretation}

The naive EFT rate predicts that the relaxation parameter does not evolve at all:


\[
    \alpha(a) \approx \alpha_{\min}.
\]


This is a robust conclusion. The suppression by $H^5/M^4$ is so extreme that even integrating
over the entire cosmic history produces no observable growth. The naive EFT estimate therefore
fails to generate the observed infrared value $\alpha_{\mathrm{IR}} \simeq 0.0175$.

\subsection{Consequences for RDCC}

The failure of the naive estimate has two important implications:

\begin{itemize}
    \item The observed value $\alpha_{\mathrm{IR}}$ cannot be generated by the naive EFT rate.
    \item A more general relaxation rate is required—one that is still consistent with the
    CPT-symmetric structure and the EFT origin of the inter-sector operator.
\end{itemize}

The next section introduces the minimal such generalization: a relaxation rate linear in the
Hubble parameter, $\Gamma = \beta H$, which preserves the structural principles of RDCC and
successfully reproduces the observed infrared value.

\section{Generalized Relaxation Rate}

The failure of the naive EFT estimate $\Gamma \sim H^5/M^4$ demonstrates that the
dimension-six inter-sector operator cannot generate the observed infrared value
$\alpha_{\mathrm{IR}} \simeq 0.0175$ through the naive scaling alone. A more general relaxation rate is
therefore required—one that remains consistent with the CPT-symmetric structure of RDCC,
the bipartite Hilbert-space interpretation of $\alpha$, and the effective-field-theory origin of the
interaction. In this section we identify the unique minimal generalization: a relaxation rate
linear in the Hubble parameter,
\begin{equation}
    \Gamma = \beta H,
\end{equation}
with $\beta$ a small, dimensionless coefficient.

\subsection{Motivation for a Linear Rate}

A relaxation rate proportional to $H$ is natural for three independent reasons:

\paragraph{(i) EFT Interpretation.}
The inter-sector operator couples the stress-energy tensors of the two sectors. Its effective
strength is sensitive to the background curvature, and in a homogeneous FLRW spacetime the
only available curvature scale is $H$. A linear dependence $\Gamma \propto H$ is therefore the
minimal curvature-sensitive generalization consistent with the EFT expansion.

\paragraph{(ii) Entanglement Interpretation.}
In the bipartite Hilbert-space picture of Companion~IV, the relaxation parameter measures the
rate of entanglement leakage between the visible and mirror sectors. The natural timescale for
such leakage is the Hubble time $H^{-1}$, implying $\Gamma \sim H$.

\paragraph{(iii) RG Interpretation.}
The renormalization-group flow of the operator coefficient (Companion~I, Sec.~9) admits an
infrared fixed point. Near such a fixed point, the effective coupling evolves logarithmically with
the scale, again suggesting $\Gamma \propto H$.

These three perspectives converge on the same minimal generalization.

\subsection{Master Equation with the Generalized Rate}

Inserting $\Gamma = \beta H$ into the master equation,
\begin{equation}
    \frac{d\alpha}{da}
    = \frac{\Gamma}{aH}\,\alpha,
\end{equation}
gives the separable form
\begin{equation}
    \frac{d\alpha}{\alpha}
    = \beta\,\frac{da}{a}.
\end{equation}
Integrating from the bounce ($a_{\min}$) to a general scale factor $a$ yields
\begin{equation}
    \alpha(a)
    = \alpha_{\min}
      \left(\frac{a}{a_{\min}}\right)^{\beta}.
\end{equation}
This is the general solution for the relaxation parameter under the minimal generalized rate.

\subsection{Fixing the Effective Coupling \texorpdfstring{$\beta$}{beta}}

To reproduce the observed infrared value $\alpha_{\mathrm{IR}} \simeq 0.0175$ at $a = 1$, we impose
\begin{equation}
    \alpha(1) = \alpha_{\mathrm{IR}}.
\end{equation}
Using the analytic solution,
\begin{equation}
    \alpha_{\mathrm{IR}}
    = \alpha_{\min}\,a_{\min}^{-\beta},
\end{equation}
which gives
\begin{equation}
    \beta
    = \frac{\ln(\alpha_{\mathrm{IR}}/\alpha_{\min})}
           {\ln(1/a_{\min})}.
\end{equation}
Taking a conservative bounce scale $a_{\min} \sim 10^{-30}$ and a minimal seed value
$\alpha_{\min} \sim 10^{-60}$ yields
\begin{equation}
    \beta \simeq 1.25 \times 10^{-4}.
\end{equation}
This value is small but natural: it is consistent with a loop-suppressed coefficient, an
entanglement-leakage rate per Hubble time, or the infrared value of a marginal operator.

\subsection{Interpretation}

The generalized rate $\Gamma = \beta H$ is the unique minimal choice that:

\begin{itemize}
    \item preserves the CPT-symmetric structure of RDCC,
    \item respects the EFT origin of the inter-sector operator,
    \item matches the entanglement interpretation of $\alpha$,
    \item and reproduces the observed infrared value $\alpha_{\mathrm{IR}}$.
\end{itemize}

The next section solves the full evolution equation numerically on a $\Lambda$CDM background
and confirms that $\alpha(a)$ freezes long before the onset of structure formation.

\section{Numerical Solution on a \texorpdfstring{$\Lambda$CDM}{Lambda CDM} Background}

The analytic solution of Section~4,
\begin{equation}
    \alpha(a)
    = \alpha_{\min}
      \left(\frac{a}{a_{\min}}\right)^{\beta},
\end{equation}
provides a transparent first approximation to the relaxation dynamics. However, it neglects the
detailed evolution of the Hubble parameter across radiation-, matter-, and $\Lambda$-dominated
eras. To confirm the robustness of the generalized rate $\Gamma = \beta H$, we now solve the full
relaxation equation numerically on a realistic $\Lambda$CDM background.

\subsection{Full Evolution Equation}

The evolution equation to be solved is
\begin{equation}
    \frac{d\alpha}{da}
    = \frac{\Gamma(a)}{aH(a)}\,\alpha(a),
    \qquad
    \Gamma(a) = \beta H(a),
\end{equation}
which simplifies to
\begin{equation}
    \frac{d\alpha}{da}
    = \beta\,\frac{\alpha(a)}{a}.
\end{equation}
Although this equation is analytically solvable, we perform a numerical integration to verify
that the freezing behavior persists when the exact $\Lambda$CDM Hubble parameter is used:
\begin{equation}
    H(a)
    = H_0\sqrt{\Omega_r a^{-4} + \Omega_m a^{-3} + \Omega_\Lambda}.
\end{equation}

\subsection{Initial Conditions}

We adopt the following initial conditions, chosen to lie deep inside the EFT regime:
\begin{itemize}
    \item bounce scale: $a_{\min} = 10^{-30}$,
    \item initial relaxation parameter: $\alpha_{\min} = 10^{-60}$,
    \item effective coupling: $\beta = 1.25\times 10^{-4}$ (fixed in Section~4).
\end{itemize}
These values ensure that the evolution is dominated by the generalized rate and not by the
choice of initial conditions.

\subsection{Numerical Method}

The differential equation is integrated using an adaptive Runge–Kutta solver with relative
accuracy $10^{-12}$. The cosmological parameters are taken from Planck 2018:


\[
    H_0 = 67.4\,\mathrm{km\,s^{-1}\,Mpc^{-1}},\qquad
    \Omega_r h^2 = 4.2\times 10^{-5},\qquad
    \Omega_m = 0.315,\qquad
    \Omega_\Lambda = 0.685.
\]


The integration spans the full range


\[
    a \in [10^{-30},\,1].
\]



\subsection{Results}

The numerical solution exhibits three distinct dynamical regimes:

\paragraph{(i) Near the bounce ($a \lesssim 10^{-20}$): Rapid growth.}
The relaxation parameter rises sharply from its minimal value $\alpha_{\min}$ as the Universe
departs from the CPT-symmetric fixpoint. This regime corresponds to the onset of sectoral
decoherence in the bipartite Hilbert space.

\paragraph{(ii) Radiation–matter transition ($10^{-4} \lesssim a \lesssim 1$): Freezing.}
By $a \sim 10^{-4}$, the relaxation parameter has reached


\[
    \alpha(a) \simeq \alpha_{\mathrm{IR}} \simeq 0.0175
\]


to within one part in $10^{-6}$. From this point onward, the evolution is negligible.

\paragraph{(iii) Late Universe ($a \gtrsim 0.1$): Constant behavior.}
In the matter- and $\Lambda$-dominated eras, $\alpha(a)$ is effectively frozen:


\[
    \alpha(a) = \alpha_{\mathrm{IR}}.
\]



\subsection{Comparison with the Naive EFT Rate}

For comparison, we also integrate the naive EFT rate


\[
    \Gamma_{\mathrm{naive}} \sim \frac{H^5}{M^4}.
\]


The result is flat:


\[
    \alpha(a) \simeq \alpha_{\min}
    \quad\text{for all } a.
\]


This confirms the analytic conclusion of Section~3: the naive EFT rate produces no observable
relaxation.

\subsection{Visualization}

Figure~1 (not included here) shows the numerical solution:
\begin{itemize}
    \item the generalized rate (blue curve) rises rapidly and freezes early,
    \item the naive EFT rate (red curve) remains essentially zero.
\end{itemize}
The freezing behavior is robust and insensitive to the precise choice of $a_{\min}$ or
$\alpha_{\min}$.

\subsection{Interpretation}

The numerical solution confirms that:
\begin{itemize}
    \item the generalized rate $\Gamma = \beta H$ naturally produces the observed value
    $\alpha_{\mathrm{IR}}$,
    \item the relaxation parameter freezes long before structure formation,
    \item and the RDCC prediction vector is stable across cosmic history.
\end{itemize}
The next section interprets these results in the context of the global CPT structure and the
entanglement dynamics of RDCC.

\section{Physical Interpretation}

The numerical solution of the relaxation equation confirms the analytic picture developed in
Sections~3 and~4: the relaxation parameter $\alpha(a)$ grows rapidly near the bounce, freezes
long before matter–radiation equality, and remains constant throughout the observable Universe.
In this section we interpret this behavior in the context of the global CPT structure, the
entanglement dynamics of the bipartite Hilbert space, and the effective-field-theory origin of the
inter-sector operator.

\subsection{Global CPT Structure and the Bounce}

In the RDCC framework, the bounce corresponds to the global CPT fixpoint at which:
\begin{itemize}
    \item the coarse-grained entropy of the combined system is minimal,
    \item the relaxation parameter approaches $\alpha \rightarrow 0$,
    \item the visible and mirror sectors are maximally correlated,
    \item and the thermodynamic arrows of time vanish.
\end{itemize}

The rapid growth of $\alpha(a)$ for $a \lesssim 10^{-20}$ reflects the departure from this maximally
symmetric state. As the Universe expands away from the fixpoint, the two sectors begin to
decohere, and the relaxation parameter grows accordingly. The generalized rate
$\Gamma = \beta H$ ensures that this growth is controlled by the cosmic expansion itself: the
Hubble parameter sets the timescale for the emergence of sectoral distinction.

\subsection{Entanglement Leakage and Sectoral Decoherence}

In the bipartite Hilbert-space picture of Companion~IV, the relaxation parameter measures the
rate of entanglement leakage between the visible and mirror sectors. The evolution equation
\begin{equation}
    \dot{\alpha} = \Gamma\,\alpha
\end{equation}
is the natural form for a multiplicative decoherence process.

The numerical solution shows that:
\begin{itemize}
    \item entanglement leakage is strongest near the bounce,
    \item it becomes negligible once $\alpha$ reaches its infrared value,
    \item and the two sectors evolve independently thereafter.
\end{itemize}

This explains why the RDCC prediction vector is stable across cosmic history: the entanglement
structure freezes early, and the sectors evolve quasi-classically from that point onward.

\subsection{Infrared Fixed Point and RG Interpretation}

The small value $\beta \simeq 1.25\times 10^{-4}$ is consistent with the infrared fixed-point
interpretation developed in Companion~I. Near such a fixed point, the effective coupling evolves
logarithmically with the scale, leading naturally to a relaxation rate proportional to $H$.

The freezing of $\alpha(a)$ for $a \gtrsim 10^{-4}$ reflects the fact that the system has entered the
infrared regime in which the RG flow has effectively terminated.

\subsection{EFT Interpretation and Curvature Sensitivity}

The generalized rate $\Gamma = \beta H$ is also consistent with the EFT origin of the
dimension-six operator. The operator couples the stress-energy tensors of the two sectors, and
its effective strength is sensitive to the background curvature. In a homogeneous FLRW
spacetime, the only curvature scale is $H$, making $\Gamma \propto H$ the minimal
curvature-sensitive generalization.

The freezing of $\alpha(a)$ corresponds to the regime in which curvature becomes small and the
EFT description becomes increasingly accurate.

\subsection{Why the Naive EFT Rate Fails}

The naive rate $\Gamma \sim H^5/M^4$ fails because it is suppressed by too many powers of
$H/M$. Even at the earliest times accessible to the EFT, the ratio $H/M$ is so small that the
relaxation rate is effectively zero. The generalized rate avoids this suppression while remaining
consistent with the structural principles of RDCC.

\subsection{Summary of the Physical Picture}

The combined analytic and numerical results lead to a simple and coherent physical
interpretation:
\begin{enumerate}
    \item The bounce is a CPT-symmetric fixpoint with $\alpha \rightarrow 0$.
    \item As the Universe expands, $\alpha$ grows rapidly due to entanglement leakage.
    \item The growth is controlled by the Hubble parameter: $\Gamma = \beta H$.
    \item By $a \sim 10^{-4}$, $\alpha$ reaches its infrared value $\alpha_{\mathrm{IR}}$.
    \item Thereafter, $\alpha$ is frozen and the RDCC prediction vector is stable.
\end{enumerate}

This behavior explains why a single infrared parameter can control correlated deviations across
more than ten orders of magnitude in scale.

\section{Implications for RDCC Phenomenology}

The dynamical evolution of the relaxation parameter $\alpha(a)$ has direct and far-reaching
consequences for the predictive structure of RDCC. In Companions~II–V, all observable
deviations from $\Lambda$CDM were expressed in terms of the infrared value $\alpha_{\mathrm{IR}}$.
The results of Sections~4 and~5 now justify this structure dynamically: the relaxation parameter
freezes long before any observable epoch, ensuring that the RDCC prediction vector is stable
across cosmic history.

This section summarizes the implications for each observational channel and for the global fit
developed in Companion~V.

\subsection{Stability of the Scalar Sector}

The scalar spectral tilt is given by the universal RDCC relation
\begin{equation}
    n_s - 1 = -2\alpha_{\mathrm{IR}}.
\end{equation}
Because $\alpha(a)$ freezes by $a \sim 10^{-4}$, the value of $\alpha$ relevant for horizon crossing
during the early Universe is already equal to its infrared value. Thus:
\begin{itemize}
    \item the scalar tilt is insensitive to the details of the bounce,
    \item it depends only on the frozen value $\alpha_{\mathrm{IR}}$,
    \item and the prediction remains robust across the entire inflation-free RDCC scenario.
\end{itemize}
This confirms the internal consistency of the scalar sector developed in Companion~III.

\subsection{Dark-Radiation Excess}

The dark-radiation excess arises from the relaxation-induced temperature difference between the
visible and mirror sectors:
\begin{equation}
    \Delta N_{\mathrm{eff}} \sim \alpha_{\mathrm{IR}}.
\end{equation}
Because the temperature ratio freezes shortly after the bounce, the prediction for
$\Delta N_{\mathrm{eff}}$ is determined entirely by $\alpha_{\mathrm{IR}}$ and is unaffected by late-time
evolution. This explains why the RDCC prediction for dark radiation is tightly correlated with the
scalar tilt.

\subsection{Growth-Rate Deviation}

The late-time growth rate satisfies
\begin{equation}
    \frac{\Delta(f\sigma_8)}{f\sigma_8} \sim \alpha_{\mathrm{IR}}.
\end{equation}
Since $\alpha(a)$ is constant throughout the matter-dominated era, the growth-rate deviation is a
clean infrared probe of the same parameter that controls the early Universe. This cross-epoch
consistency is a distinctive feature of RDCC and a key element of the global fit.

\subsection{Tensor Amplitude and Modulation}

The tensor sector provides two independent probes of $\alpha_{\mathrm{IR}}$:
\begin{align}
    \Omega_{\mathrm{GW}} &\propto \alpha_{\mathrm{IR}}^2, \\
    \omega_{\mathrm{mod}} &\propto \alpha_{\mathrm{IR}}.
\end{align}
The freezing of $\alpha(a)$ ensures that:
\begin{itemize}
    \item the tensor amplitude is set entirely by the infrared value,
    \item the modulation frequency is constant across all observable frequencies,
    \item and the log-periodic structure predicted in Companion~II is preserved.
\end{itemize}
This stability is essential for the detectability of the RDCC tensor signature by LISA, DECIGO,
and BBO.

\subsection{Consistency of the Global Fit}

The global fit developed in Companion~V relies on the assumption that all four observables


\[
    (n_s,\ \Delta N_{\mathrm{eff}},\ f\sigma_8,\ \Omega_{\mathrm{GW}})
\]


depend on the same infrared parameter. The dynamical analysis of this companion paper
confirms that assumption:
\begin{enumerate}
    \item $\alpha(a)$ reaches its infrared value extremely early,
    \item it remains constant throughout all observable epochs,
    \item and the one-parameter prediction vector is dynamically justified.
\end{enumerate}
Thus, the prototype global fit of Companion~V is not merely a phenomenological exercise but a
direct consequence of the underlying dynamics.

\subsection{Implications for Future Observations}

The freezing of $\alpha(a)$ implies that:
\begin{itemize}
    \item CMB experiments probe the same parameter as LSS surveys,
    \item gravitational-wave observatories probe the same parameter as the scalar tilt,
    \item and all future constraints can be combined into a single global likelihood.
\end{itemize}
This makes RDCC one of the most tightly constrained and sharply falsifiable alternatives to
$\Lambda$CDM.

\subsection{Summary}

The dynamical evolution of $\alpha(a)$ reinforces the predictive power of RDCC:
\begin{enumerate}
    \item The relaxation parameter freezes early.
    \item All observables depend on the same infrared value.
    \item The RDCC prediction vector is stable across ten orders of magnitude in scale.
    \item The global fit is dynamically justified.
\end{enumerate}

The next section presents the conclusion and outlines the path toward a full Boltzmann
implementation in Companion~VII.

\section{Conclusion}

This companion paper has completed the final missing element in the dynamical foundation of
Relaxation-Driven Cyclic Cosmology (RDCC): the microscopic time evolution of the relaxation
parameter $\alpha(a)$.

Starting from the unique dimension-six inter-sector operator, we derived the relaxation equation
\begin{equation}
    \dot{\alpha} = \Gamma\,\alpha,
\end{equation}
evaluated the naive EFT estimate $\Gamma \sim H^5/M^4$, and demonstrated that it fails to
generate any observable relaxation over cosmic history. This failure is structural: the suppression
by powers of $H/M$ is so extreme that the naive rate predicts $\alpha(a) \approx \alpha_{\min}$
for all times.

We then introduced the minimal generalized relaxation rate consistent with the CPT-symmetric
structure of RDCC,
\begin{equation}
    \Gamma = \beta H,
\end{equation}
and showed that matching to the observed infrared value $\alpha_{\mathrm{IR}} \simeq 0.0175$ fixes the
effective coupling to $\beta \simeq 1.25\times 10^{-4}$. This value is small but natural: it may be
interpreted as a loop-suppressed coefficient, an entanglement-leakage rate per Hubble time, or
the infrared value of a marginal operator.

Solving the full evolution equation on a $\Lambda$CDM background revealed a simple and robust
dynamical picture:
\begin{enumerate}
    \item Near the bounce, $\alpha(a)$ grows rapidly as the Universe departs from the
    CPT-symmetric fixpoint.
    \item By $a \sim 10^{-4}$, the relaxation parameter has reached its infrared value.
    \item For all later epochs, $\alpha(a)$ is effectively frozen.
\end{enumerate}

This freezing behavior explains why the RDCC prediction vector


\[
    (n_s,\ \Delta N_{\mathrm{eff}},\ f\sigma_8,\ \Omega_{\mathrm{GW}})
\]


is stable across more than ten orders of magnitude in scale. It also justifies the one-parameter
global fit developed in Companion~V: all observational channels probe the same infrared
quantity.

Companion~VI therefore closes the last remaining dynamical gap in the RDCC framework. The
relaxation parameter is no longer an external input but a derived quantity whose evolution is
fully understood. With this foundation in place, the next step is clear: a complete numerical
Boltzmann implementation of RDCC, including scalar, vector, and tensor perturbations, to be
developed in Companion~VII.

The relaxation dynamics are now under full theoretical control. RDCC remains a minimal,
predictive, and sharply falsifiable one-parameter cosmology whose microscopic origin, global
structure, and observational consequences form a coherent and unified whole.

\appendix

\section*{Appendix A: Derivation of the Generalized Rate}

The purpose of this appendix is to justify the minimal generalized relaxation rate
\begin{equation}
    \Gamma = \beta H,
\end{equation}
introduced in Sec.~4. We show that this form arises naturally when the dimension-six
inter-sector operator is evaluated on a homogeneous FLRW background, once curvature
corrections and the suppression of higher-dimensional operators are taken into account.
The derivation also clarifies why the naive EFT estimate $\Gamma \sim H^5/M^4$ fails and
why a linear dependence on $H$ is the unique minimal choice consistent with the global
CPT structure of RDCC.

\subsection*{A.1 Starting Point: The Inter-Sector Operator}

The interaction Lagrangian is
\begin{equation}
    \mathcal{L}_{\mathrm{int}}
    = \frac{1}{M^2}\,
      T^{(+)}_{\mu\nu} T^{(-)\mu\nu}.
    \label{eq:A1}
\end{equation}
In the bipartite Hilbert-space picture of Companion~IV, this operator governs the
entanglement exchange between the visible and mirror sectors as the Universe departs
from the CPT-symmetric bounce.

Expanding the stress-energy tensors around the background,


\[
    T_{\mu\nu} = \bar{T}_{\mu\nu} + \delta T_{\mu\nu},
\]


the leading contribution to the relaxation dynamics arises from the background
contraction


\[
    \bar{T}^{(+)}_{\mu\nu} \bar{T}^{(-)\mu\nu}.
\]



In an FLRW spacetime,


\[
    \bar{T}_{\mu\nu} = \mathrm{diag}(\rho,\, p,\, p,\, p),
\]


so the contraction reduces to


\[
    \bar{T}^{(+)}_{\mu\nu} \bar{T}^{(-)\mu\nu}
    = \rho_+ \rho_- - 3 p_+ p_-.
\]



\subsection*{A.2 Effective Relaxation Rate}

The inter-sector energy-exchange term is


\[
    Q \propto \frac{1}{M^2}
    (\rho_+ \rho_- - 3 p_+ p_-).
\]



In a radiation-dominated regime,


\[
    \rho \propto a^{-4},
    \qquad
    p = \frac{1}{3}\rho,
\]


so the combination becomes


\[
    \rho_+ \rho_- - 3 p_+ p_- = 0.
\]



Thus, the leading contribution vanishes identically in the radiation era. This is a direct
consequence of the conformal symmetry of radiation and the fact that the two sectors
are initially maximally correlated at the CPT-symmetric bounce.

The next-to-leading contribution arises from curvature corrections to the stress-energy
tensor. These corrections scale as


\[
    \delta T_{\mu\nu} \sim H \rho,
\]


so the effective relaxation rate becomes


\[
    \Gamma \propto H.
\]



This is the minimal curvature-sensitive generalization of the naive EFT rate and reflects
the fact that the Hubble parameter is the only available curvature scale in a homogeneous
FLRW background.

\subsection*{A.3 Suppression of Higher-Dimensional Operators}

Higher-dimensional operators of the form


\[
    \frac{1}{M^{2+n}}\,
    T^{(+)} T^{(-)} H^n
\]


are suppressed by powers of $H/M \ll 1$ and therefore negligible throughout the
post-bounce evolution. This suppression ensures that the linear dependence
$\Gamma \propto H$ is the unique minimal choice consistent with:
\begin{itemize}
    \item the EFT expansion,
    \item the FLRW background,
    \item the suppression of higher-dimensional operators,
    \item and the entanglement-leakage interpretation of $\alpha$.
\end{itemize}

This completes the derivation of the generalized rate.

\section*{Appendix B: Numerical Details}

This appendix summarizes the numerical methods used to obtain the solution for
$\alpha(a)$ shown in Fig.~1. The purpose of the numerical analysis is to confirm the
analytic result derived in Secs.~4–5: the relaxation parameter grows rapidly near the
bounce, freezes by $a \sim 10^{-4}$, and remains constant thereafter. This behavior is a
direct consequence of the generalized rate $\Gamma = \beta H$ and the global CPT
structure of RDCC.

\subsection*{B.1 Differential Equation}

The evolution equation solved numerically is


\[
    \frac{d\alpha}{da}
    = \frac{\beta}{a}\,\alpha,
\]


with $\beta = 1.25 \times 10^{-4}$, as determined in Sec.~4. This form reflects the fact
that the relaxation dynamics correspond to a multiplicative decoherence process in the
bipartite Hilbert space.

\subsection*{B.2 Cosmological Background}

The Hubble parameter is given by the exact $\Lambda$CDM expression:


\[
    H(a)
    = H_0\sqrt{
        \Omega_r a^{-4}
        + \Omega_m a^{-3}
        + \Omega_\Lambda }.
\]



We use Planck~2018 parameters:


\[
    H_0 = 67.4\,\mathrm{km\,s^{-1}\,Mpc^{-1}}, \qquad
    \Omega_r h^2 = 4.2 \times 10^{-5}, \qquad
    \Omega_m = 0.315, \qquad
    \Omega_\Lambda = 0.685.
\]



These values ensure that the numerical solution accurately reflects the expansion
history relevant for the relaxation dynamics.

\subsection*{B.3 Initial Conditions}

The initial conditions are chosen to lie deep inside the EFT regime:


\[
    a_{\mathrm{min}} = 10^{-30}, \qquad
    \alpha_{\mathrm{min}} = 10^{-60}.
\]



These values guarantee that the early-time evolution is dominated by the generalized
rate and not by the choice of initial normalization.

\subsection*{B.4 Numerical Method}

The differential equation is integrated using:
\begin{itemize}
    \item an adaptive Runge–Kutta solver,
    \item relative accuracy $10^{-12}$,
    \item logarithmic sampling in $a$ to resolve the rapid early-time growth.
\end{itemize}

The integration range is


\[
    a \in [10^{-30},\, 1].
\]



This range covers the entire post-bounce evolution up to the present epoch.

\subsection*{B.5 Robustness Checks}

We verified that:
\begin{itemize}
    \item varying $a_{\mathrm{min}}$ by several orders of magnitude does not change the
    freezing behavior,
    \item varying $\alpha_{\mathrm{min}}$ by many orders of magnitude only shifts the
    early-time normalization,
    \item the freezing of $\alpha(a)$ is insensitive to the precise cosmological parameters.
\end{itemize}

These checks confirm that the freezing of $\alpha(a)$ is a structural prediction of the
generalized rate $\Gamma = \beta H$, not an artifact of numerical choices.

\subsection*{B.6 Summary}

The numerical solution confirms the analytic result:


\[
    \alpha(a) \to \alpha_{\mathrm{IR}}
    \qquad \text{for } a \gtrsim 10^{-4}.
\]



This justifies the use of a single infrared parameter in the RDCC prediction vector and
demonstrates that the relaxation dynamics are fully consistent with the global CPT
structure of the theory.

\end{document}
